% !TeX program = lualatex
% ================================================================
% Urdu Parallel Text Version of TrigBearingsStarter01
%
% Generated by the server using the following settings:
%
%   language            Urdu (urd)
%   text direction      right to left
%   font                Noto Nastaliq Urdu
%   fallback font       Noto Serif
%   built from          latex/starters/targeted/KS4/geometry/trigBearingsStarter01.tex
%   vocabulary key      latex/starters/targeted/KS4/geometry/trigBearingsStarter01/L2/urd/trigBearingsStarter01_urduKey.tex
%   tier 3 vocabulary   green   RGB 0 128 0
%   English text        black   RGB 0 0 0
%   Urdu text           purple  RGB 112 48 160
%   layout macros       version 3
%
% Please don't edit any information in this comment block - the server
% compares it against the current settings and rebuilds this file when
% they differ, so changes here would be overwritten. However, feel free
% to make updates to the LaTeX below if you want to edit it.
% ================================================================

\documentclass[aspectratio=169]{beamer}
% ================================================================
% Copied from the English sheet this was translated from, so that
% anything it sets up for itself still works here
% ================================================================
\usepackage{amsmath}
\usepackage{tikz}
\usetikzlibrary{angles,quotes,calc}

% ================================================================
% Answer-overlay helpers
% ================================================================
\newcommand{\ablank}[1]{%
  \alt<2>{\textcolor{red}{#1}}{\underline{\phantom{#1}}}%
}
\newcommand{\ashow}[1]{\uncover<2->{\textcolor{red}{\small #1}}}
\newcommand{\ashowq}[1]{\alt<2->{\textcolor{red}{\small #1}}{?}}

% Answers drawn straight on to a TikZ diagram, revealed on overlay 2 like \ashow.
% Space is reserved on overlay 1, so the diagram does not move between slides.
\usetikzlibrary{overlay-beamer-styles}
\tikzset{
  ans/.style     = {red, thick, visible on=<2->},
  ansfill/.style = {red!15, visible on=<2->},
  anslab/.style  = {red, font=\tiny, visible on=<2->},
}

\title{Starter}
\author{}
\date{}

\setbeamertemplate{navigation symbols}{}

% Characters this language's font has not got are borrowed from another,
% rather than being silently dropped - see docs/translated-sheet-latex.md
\directlua{%
  if luaotfload and luaotfload.add_fallback then
    luaotfload.add_fallback("eallatinfallback",
      { "Noto Serif:mode=harf;" })
  else
    tex.error("this luaotfload is too old to register a fallback font")
  end
}
\usepackage[english,bidi=basic]{babel}
\babelprovide[import]{urdu}
\babelfont[urdu]{rm}[Renderer=Harfbuzz, BoldFont={Noto Nastaliq Urdu}, ItalicFont={Noto Nastaliq Urdu}, BoldItalicFont={Noto Nastaliq Urdu}, RawFeature={fallback=eallatinfallback}]{Noto Nastaliq Urdu}
\babelfont[urdu]{sf}[Renderer=Harfbuzz, BoldFont={Noto Nastaliq Urdu}, ItalicFont={Noto Nastaliq Urdu}, BoldItalicFont={Noto Nastaliq Urdu}, RawFeature={fallback=eallatinfallback}]{Noto Nastaliq Urdu}
\newcommand{\ealtext}[1]{\foreignlanguage{urdu}{#1}}
\newcommand{\ealtextblock}[1]{{\raggedleft\foreignlanguage{urdu}{#1}\par}}
% ================================================================
% EAL layout helpers
% ================================================================
\usepackage{xcolor}

\definecolor{ealkeycolour}{RGB}{0,128,0}
\definecolor{ealenglish}{RGB}{0,0,0}
\definecolor{ealtwocolour}{RGB}{112,48,160}

\newsavebox{\ealboxen}
\newsavebox{\ealboxtr}
\newlength{\ealglosswd}
\newlength{\ealglossraise}
\setlength{\ealglossraise}{1.15em}

% a tier 3 word, in English and in translation alike
\newcommand{\ealkey}[1]{\textcolor{ealkeycolour}{#1}}
\newcommand{\ealkeytr}[1]{\textcolor{ealkeycolour}{\ealtext{#1}}}

% One translated word above one English word. Built from kernel
% primitives, never a tabular, which would break wherever the sheet
% loads array, booktabs or tabularx. The box is as wide as the wider of
% the two, so a long translation pushes its neighbours aside instead of
% overlapping them, and it claims its full height so a table row or a
% TikZ node makes room for it.
\newcommand{\ealgloss}[2]{%
  \sbox{\ealboxen}{\ealkey{#1}}%
  \sbox{\ealboxtr}{\small\textcolor{ealtwocolour}{\ealtext{#2}}}%
  \setlength{\ealglosswd}{\wd\ealboxen}%
  \ifdim\wd\ealboxtr>\ealglosswd\setlength{\ealglosswd}{\wd\ealboxtr}\fi
  \leavevmode
  \makebox[\ealglosswd][c]{%
    \makebox[0pt][c]{\usebox{\ealboxen}}%
    \makebox[0pt][c]{\raisebox{\ealglossraise}{\usebox{\ealboxtr}}}%
  }%
}

% opens up line spacing around a block of glosses, so the raised
% translations sit clear of the line above
\newenvironment{ealglossed}{\par\linespread{2.1}\selectfont\ignorespaces}{\par}

% a whole translated sentence above its English counterpart. block level,
% so it does not belong inside a tabular cell
\newcommand{\ealpara}[2]{%
  \par\addvspace{0.5\baselineskip}%
  {\color{ealtwocolour}\ealtextblock{#1}}%
  \penalty200
  {\color{ealenglish}#2\par}%
  \addvspace{0.5\baselineskip}%
}

\begin{document}

\begin{frame}[allowframebreaks]
  \frametitle{Starter}

  \begin{columns}[t,onlytextwidth]

    % ===== LEFT COLUMN =====
    \begin{column}{0.4\textwidth}
      % ------------------ Q1 ------------------
      \textbf{1.}

      \vspace{0.7em}
      \ealpara{شمال، جنوب، مشرق اور مغرب میں سے کون سے \ealkeytr{بنیادی سمتوں} کے جوڑے \ealkeytr{قائمہ زاویہ} پر ملتے ہیں؟}{Which pairs of \ealkey{cardinal directions} (North, South, East, West) meet at \ealkey{right angles}?}

      \ealpara{شمال اور مشرق، شمال اور مغرب، جنوب اور مشرق، جنوب اور مغرب۔}{North and East, North and West, South and East, South and West.}
      \ashow{North and East, North and West, South and East, South and West.}

      % ----- Flexible space -----
      \vfill
      \vspace{3em}

      % ------------------ Q3 ------------------
      \textbf{3.}

      \begin{tikzpicture}[scale=0.2, thick]
        % Right triangle (hypotenuse is 13)
        \coordinate (A) at (0,0);
        \coordinate (B) at (5,0);
        \coordinate (C) at (5,13);

        % Triangle
        \draw (A) -- (B) -- (C) -- cycle;

        % Right angle marker
        \draw ($(B)+(0,1)$) -- ($(B)+(-1,1)$) -- ($(B)+(-1,0)$);

        % Side lengths
        \node[below] at ($(A)!0.5!(B)$) {$5$ cm};
        \node[right] at ($(B)!0.5!(C)$) {$x$};
        \node[above left] at ($(A)!0.5!(C)$) {$13$ cm};

      \end{tikzpicture}

      \vspace{0.5em}
      \ealpara{گمشدہ ضلع معلوم کریں۔}{Find the missing side.}

      \ashow{$x = 12$ cm}

    \end{column}

    % ===== RIGHT COLUMN =====
    \begin{column}{0.6\textwidth}
      % ------------------ Q2 ------------------
      
      \textbf{2.}

      \begin{columns}[t,onlytextwidth]
      \begin{column}{0.5\textwidth}

      \begin{tikzpicture}[scale=0.8, thick]

        % Coordinates of triangle (hypotenuse longest)
        \coordinate (A) at (0,0);
        \coordinate (B) at (1.8,0);
        \coordinate (C) at (0,2.2);

        % Draw triangle
        \draw (A) -- (B) -- (C) -- cycle;

        % Right angle marker
        \draw ($(A)+(0,0.4)$) -- ($(A)+(0.4,0.4)$) -- ($(A)+(0.4,0)$);

        % Label sides
        \node[below] at ($(A)!0.5!(B)$) {3 cm};
        \node[left]  at ($(A)!0.5!(C)$) {4 cm};
        \node[above right] at ($(B)!0.5!(C)$) {5 cm};

      \end{tikzpicture}
\end{column}
      \begin{column}{0.5\textwidth}
      \ealpara{کون سی ضلع \ealkeytr{وتر} ہے؟}{Which side is the \ealkey{hypotenuse}?}

      \ealpara{وہ ضلع جس پر 5 cm لکھا ہے وہ \ealkeytr{وتر} ہے۔}{The side labelled 5 cm is the \ealkey{hypotenuse}.}
      \ashow{The side labelled 5 cm is the \ealkey{hypotenuse}.}
\end{column}

    \end{columns}

\textbf{4.}

\begin{columns}[t,onlytextwidth]
  % ===== Diagram Column =====
  \begin{column}{0.45\textwidth}

    \begin{tikzpicture}[scale=0.8, thick]
      % Triangle for inverse sine
      \coordinate (A) at (0,0);
      \coordinate (B) at (3,0);
      \coordinate (C) at (3,2);

      % Draw triangle
      \draw (A) -- (B) -- (C) -- cycle;

      % Right angle
      \draw ($(B)+(0,0.4)$) -- ($(B)+(-0.4,0.4)$) -- ($(B)+(-0.4,0)$);

      % Side labels
      \node[right] at ($(B)!0.5!(C)$) {$2.5$ m};  % opposite
      \node[above left] at ($(A)!0.5!(C)$) {$4$ m}; % hypotenuse

      % Angle θ at A
      \pic[draw,angle radius=10mm,"$\theta$"{scale=1}] {angle = B--A--C};

    \end{tikzpicture}

  \end{column}

  % ===== Options Column =====
  \begin{column}{0.55\textwidth}

    \ealpara{$\theta$ معلوم کریں۔}{Find $\theta$.}

     \fontsize{10}{12}\selectfont

    \begin{itemize}
        \item[(A)] $\displaystyle \theta = \sin^{-1}\!\left(\frac{2.5}{4}\right)$
        \item[(B)] $\displaystyle \theta = \cos^{-1}\!\left(\frac{2.5}{4}\right)$
        \item[(C)] $\displaystyle \theta = \tan^{-1}\!\left(\frac{2.5}{4}\right)$
    \end{itemize}

    \ealpara{جواب: (A)}{Answer: (A)}
    \ashow{Answer: (A)}

  \end{column}
\end{columns}
\end{column}
\end{columns}

\end{frame}

\end{document}